\documentclass[11pt]{article}

%==================== packages ====================
\usepackage[margin=1in]{geometry}
\usepackage{amsmath,amssymb,amsthm,mathtools}
\usepackage{enumitem}
\usepackage{xcolor}
\usepackage{hyperref}
\hypersetup{colorlinks=true,linkcolor=blue!50!black,citecolor=blue!50!black}

%==================== macros ====================
\newcommand{\cP}{\mathcal{P}}
\newcommand{\cQ}{\mathcal{Q}}
\newcommand{\Z}{\mathbb{Z}}
\newcommand{\F}{\mathbb{F}}
\newcommand{\bG}{\mathbb{G}}
\newcommand{\bP}{\mathbb{P}}
\newcommand{\bA}{\mathbb{A}}
\newcommand{\bW}{\mathbb{W}}
\newcommand{\Oo}{\mathcal{O}}
\newcommand{\mdls}{\mathfrak{m}}
\newcommand{\ndls}{\mathfrak{n}}
\DeclareMathOperator{\Frac}{Frac}
\DeclareMathOperator{\Gal}{Gal}
\DeclareMathOperator{\Spec}{Spec}
\newcommand{\AS}{\wp}
\newcommand{\vf}{\vec f}
\newcommand{\va}{\vec a}
\newcommand{\vX}{\vec X}
\newcommand{\vY}{\vec Y}
\newcommand{\vZ}{\vec Z}

\theoremstyle{plain}
\newtheorem{theorem}{Theorem}
\newtheorem{lemma}[theorem]{Lemma}
\newtheorem{prop}[theorem]{Proposition}
\newtheorem{cor}[theorem]{Corollary}
\theoremstyle{definition}
\newtheorem{definition}[theorem]{Definition}
\newtheorem{remark}[theorem]{Remark}

\title{Closing the gaps:\\ a complete proof that $\cP^{[\deg\ndls]}$ is unramified at $\eta_\ndls$ for $G=\Z/p^r\Z$}
\author{Working note}
\date{}

\begin{document}
\maketitle

\begin{abstract}
The note \emph{Completing the induction} proves the theorem
``$\cP^{[\deg\ndls]}$ is unramified at $\eta_\ndls$ for $G=\Z/p^r\Z$'' \emph{conditionally}
on two unproven steps (Gaps~1 and~2 of the review): (1) that the Witt datum of the torsor
at $p_\infty$ may be taken to be a \emph{reduced vector of polynomials in $x^{-1}$}
satisfying the pole bound $p^{r-1-j}m_j<d$, and (2) that the symmetric-power cover at
$\eta_C$ is the Artin--Schreier--Witt cover of the Witt sum
$\bigoplus_{i=1}^d\vf(x_i)$. This note supplies both missing results.
Gap~1 is closed by a $W_r$-analogue of the thesis's realization lemma over $\bG_m$:
every class in $W_r(k((x)))/\AS$ admits a reduced representative with all components in
$x^{-1}k[x^{-1}]$, produced by induction on $r$ through the truncation
$W_r\twoheadrightarrow W_{r-1}$, using that adding $V^{r-1}(h)$ is carry-free. (The
``principal part'' step of the earlier note is \emph{not} used: componentwise truncation of
a Witt vector is not a legitimate operation.) Gap~2 is closed by the Witt-coordinate
analogue of the thesis's invariant-ring computation: the Witt sum
$\vY=\bigoplus_i\vX_i$ is invariant under the twisted $G^{d-1}$-action and under $S_d$,
satisfies $\AS\vY=\bigoplus_i\vf(x_i)$, and the resulting map of $G$-torsors is
automatically an isomorphism. Together with the (verified) degree count of the earlier
note, reproduced here for completeness, this yields an unconditional proof of the theorem,
directly for all $r$ at once. The single external input remains the
Artin--Schreier--Witt break formula, stated here in the precise form and direction used.
\end{abstract}

\tableofcontents

%====================================================================
\section{Setting, conventions, and what is being proved}
%====================================================================

Throughout, $k$ is an algebraically closed field of characteristic $p>0$ and
$K=k((x))$. We prove:

\begin{theorem}[{main theorem; \texttt{theorem:ramification\_after\_blowup\_is\_tame\_for\_p\_power}}]\label{thm:main}
Let $C$ be a projective smooth geometrically connected curve over $k$, let
$G=\Z/p^r\Z$, and let $\cP\to U=C\setminus\mdls$ be a $G$-torsor with ramification bounded
by $\mdls$. Let $\ndls=nP\subset\mdls$ and $d:=\deg\ndls=n$. Then $\cP^{[d]}$ is
unramified at $\eta_\ndls$.
\end{theorem}

Here, as in the thesis, ``ramification bounded by $d$ at $P$'' means that the $d$-th
\emph{upper-numbering} ramification group of the local extension at $P$ is trivial
(\texttt{definition:local\_ramification\_boundness}); equivalently, the largest
upper-numbering break $u$ of that extension satisfies $u<d$ \emph{strictly}.

\paragraph{Standing reduction.}
As in the current (partial) proof in the thesis, we may assume that $k$ is algebraically
closed and that $\cP\to U$ is \emph{totally ramified} at $P$ with inertia equal to all of
$G$: if the image $H$ of inertia at $P$ is a proper subgroup, decompose
$\cP\to\cP_{G/H}\to U$; then $\cP_{G/H}$ is unramified at $P$, extends over $P$, and is
locally trivial at any point over $P$ (as $k=\bar k$), which replaces $(G,\cP)$ by
$(H,\cP\to\cP_{G/H})$ with $H\cong\Z/p^s\Z$. This reduction is valid for all $r$ and is
kept unchanged; we do not repeat it below. After it, the restriction
$\cQ:=\cP_x$ of $\cP$ to the punctured formal disc $\Spec k((x))$ at $P$ (with $x$ a
uniformizer at $P$) is a \emph{connected} $\Z/p^r\Z$-torsor, i.e.\ $\cQ=\Spec L$ for a
cyclic, totally ramified extension $L/K$ of degree $p^r$, with ramification bounded by
$d$.

\paragraph{Witt vector conventions.}
$W_r$ denotes $p$-typical Witt vectors of length $r$, with components indexed
$0,\dots,r-1$; $\oplus,\ominus$ denote Witt addition and subtraction, given by the
universal polynomials $S_n$: the $n$-th component of $\va\oplus\vec b$ is
\begin{equation}\label{eq:addition-shape}
(\va\oplus\vec b)_n \;=\; a_n+b_n+C_n(a_0,\dots,a_{n-1};b_0,\dots,b_{n-1})
\end{equation}
for a universal polynomial $C_n$ over $\Z$ in the \emph{lower} components only.
On $\F_p$-algebras the Frobenius $F$ is componentwise:
$F(a_0,\dots,a_{r-1})=(a_0^p,\dots,a_{r-1}^p)$; we set
$\AS:=F-\mathrm{id}$ (computed with $\ominus$), so for $r=1$, $\AS(u)=u^p-u$.
The Verschiebung is $V(a_0,\dots,a_{r-2})=(0,a_0,\dots,a_{r-2})$, and
$t\colon W_r\to W_{r-1}$, $t(a_0,\dots,a_{r-1})=(a_0,\dots,a_{r-2})$, is the truncation.
$W_r$ is a functor: a ring homomorphism $\sigma\colon A\to B$ induces
$W_r(\sigma)\colon W_r(A)\to W_r(B)$ acting \emph{componentwise},
$W_r(\sigma)(a_0,\dots,a_{r-1})=(\sigma a_0,\dots,\sigma a_{r-1})$, and this is a ring
homomorphism commuting with $F$, $V$, $t$, hence with $\AS$. In particular a ring
automorphism $\sigma$ of $A$ commutes with $\oplus$ and $\ominus$, and a Witt vector
$\va\in W_r(A)$ satisfies $W_r(\sigma)(\va)=\va$ if and only if each component $a_n$ is
$\sigma$-invariant.
By Artin--Schreier--Witt (ASW) theory, over a field $M\supseteq\F_p$ one has
$H^1(M,\Z/p^r\Z)\cong W_r(M)/\AS W_r(M)$, the class of $\vec g$ corresponding to the
torsor $\Spec M[\vX]/(\AS\vX\ominus\vec g)$ (notation of
Lemma~\ref{lem:asw-etale} below); connected torsors correspond to classes of order
$p^r$, and the extension is then cyclic of degree $p^r$, totally ramified when
$M=k((x))$ with $k$ algebraically closed.

We record the elementary Witt bookkeeping used repeatedly.

\begin{lemma}\label{lem:bookkeeping}
Let $A$ be an $\F_p$-algebra.
\begin{enumerate}[label=(\roman*)]
  \item $t\colon W_r(A)\to W_{r-1}(A)$ is a ring homomorphism commuting with $F$, hence
        with $\AS$; its kernel is $V^{r-1}A=\{(0,\dots,0,c):c\in A\}$.
  \item $V^{r-1}\colon A\to W_r(A)$ is additive, and $\AS\circ V^{r-1}=V^{r-1}\circ\AS$.
  \item (carry-free top-level addition) For all $\va\in W_r(A)$ and $c\in A$,
        \[
        \va\oplus V^{r-1}(c)=(a_0,\dots,a_{r-2},\,a_{r-1}+c).
        \]
\end{enumerate}
\end{lemma}

\begin{proof}
(i) is standard ($t$ is a natural transformation of ring functors; $F$ is componentwise
in characteristic $p$), and the kernel is computed componentwise.
(ii) Additivity of $V$ is standard; since $F$ is componentwise it commutes with $V$ on
$\F_p$-algebras, hence so does $\AS=F-\mathrm{id}$.
(iii) This is a universal identity, so it suffices to prove it over
$\Z[a_0,\dots,a_{r-1},c]$, where the ghost components
$w_n(\va)=\sum_{j\le n}p^ja_j^{p^{n-j}}$ are injective jointly with torsion-freeness.
Put $\vec b=(0,\dots,0,c)$ and $\vec s=\va\oplus\vec b$. For $n<r-1$ we have
$w_n(\vec b)=0$, so $w_n(\vec s)=w_n(\va)$, and by induction on $n$ (solving for $s_n$
using torsion-freeness) $s_n=a_n$ for $n\le r-2$. For $n=r-1$,
$w_{r-1}(\vec b)=p^{r-1}c$ gives
$p^{r-1}s_{r-1}=p^{r-1}a_{r-1}+p^{r-1}c$, i.e.\ $s_{r-1}=a_{r-1}+c$.
\end{proof}

\begin{lemma}[ASW covers are \'etale torsors]\label{lem:asw-etale}
Let $A$ be an $\F_p$-algebra and $\vec g\in W_r(A)$. Let
\[
B \;:=\; A[X_0,\dots,X_{r-1}]\big/\big(\text{the $r$ components of } \AS(\vX)\ominus\vec g\big),
\qquad \vX=(X_0,\dots,X_{r-1}).
\]
Then $B$ is finite free over $A$ with basis
$\{\prod_{n}X_n^{a_n}:0\le a_n\le p-1\}$, and \'etale over $A$. The translation action
of $W_r(\F_p)\cong\Z/p^r\Z$, $\vX\mapsto\vX\oplus\underline m$ (where $\underline m$ is
the image of $m$ under $\Z/p^r\Z\cong W_r(\F_p)$), makes $\Spec B$ a $\Z/p^r\Z$-torsor
over $\Spec A$.
\end{lemma}

\begin{proof}
By \eqref{eq:addition-shape} and componentwise $F$, the $n$-th defining equation has the
triangular form
\[
X_n^p - X_n - g_n + C_n'(X_0,\dots,X_{n-1};\,g_0,\dots,g_{n-1}) \;=\; 0
\]
for universal polynomials $C_n'$. Triangularity gives the free basis; the Jacobian
matrix of the system is lower triangular with diagonal entries
$pX_n^{p-1}-1=-1$, a unit, so $B$ is \'etale over $A$. Finally, $\Spec B$ is the
pullback along $\vec g\colon\Spec A\to\bW_r$ of the Lang covering
$\AS\colon\bW_r\to\bW_r$ of the Witt group scheme over $\F_p$, which is a surjective
\'etale group homomorphism with kernel $W_r(\F_p)$; hence $\Spec B$ is a
$W_r(\F_p)$-torsor.
\end{proof}

\begin{lemma}\label{lem:torsor-mor}
Let $G$ be a finite group scheme over a base $T$ and let $\phi\colon\cP_1\to\cP_2$ be a
$G$-equivariant $T$-morphism of $G$-torsors (for the fppf or \'etale topology). Then
$\phi$ is an isomorphism.
\end{lemma}

\begin{proof}
Locally on a cover trivializing both torsors, $\phi$ is a $G$-equivariant endomorphism of
the trivial torsor $G_T$, i.e.\ right translation by a section of $G$, hence an
isomorphism; being an isomorphism descends.
\end{proof}

%====================================================================
\section{Gap 1: reduced polynomial representatives and realization over $\bG_m$}
%====================================================================

The earlier note took $\vf$ to be ``the principal part'' of the Witt datum. As the review
observes, this is not a legitimate operation (Witt addition is not componentwise, so
componentwise truncation to polar parts changes the ASW class), and in any case no
$W_r$-realization statement was available in the thesis, whose realization lemma
(\texttt{lemma:realization\_over\_Gm}) covers only the tame case and $G=\Z/p\Z$. This
section proves the correct statement.

\begin{definition}\label{def:reduced}
For $g\in k((x))$ put $m(g):=\max\bigl(0,-v_x(g)\bigr)$ (the pole order, or $0$ if $g$ is
regular). A vector $\vec g=(g_0,\dots,g_{r-1})\in W_r(k((x)))$ is \emph{reduced} if for
every $j$, either $m(g_j)=0$ or $p\nmid m(g_j)$. It is a \emph{reduced polynomial vector}
if moreover every $g_j\in x^{-1}k[x^{-1}]$ (so that $m(g_j)=\deg_{x^{-1}}g_j$, and
$m(g_j)=0$ if and only if $g_j=0$).
\end{definition}

\begin{lemma}[surjectivity of $\AS$ on integral Witt vectors]\label{lem:wp-surj}
$\AS\colon W_r(k[[x]])\to W_r(k[[x]])$ is surjective.
\end{lemma}

\begin{proof}
Induction on $r$. For $r=1$ this is proved in the thesis
(\texttt{lemma:realization\_over\_Gm}, first bullet): given
$h=\sum_{i\ge0}b_ix^i$, solve $u^p-u=h$ coefficientwise, using that $k$ is algebraically
closed for the constant term. For $r>1$, let $\vec h\in W_r(k[[x]])$. By induction choose
$\vec u'\in W_{r-1}(k[[x]])$ with $\AS\vec u'=t(\vec h)$, and lift it to
$\tilde u:=(\vec u',0)\in W_r(k[[x]])$. By Lemma~\ref{lem:bookkeeping}(i),
$t\bigl(\AS\tilde u\ominus\vec h\bigr)=\AS\vec u'\ominus t(\vec h)=0$, so
$\AS\tilde u\ominus\vec h=V^{r-1}(g)$ for some $g\in k[[x]]$. Choose $w\in k[[x]]$ with
$\AS w=-g$ (case $r=1$ again). Then by Lemma~\ref{lem:bookkeeping}(ii) and additivity of
$V^{r-1}$,
\[
\AS\bigl(\tilde u\oplus V^{r-1}(w)\bigr)
=\AS\tilde u\oplus V^{r-1}(\AS w)
=\bigl(\vec h\oplus V^{r-1}(g)\bigr)\oplus V^{r-1}(-g)=\vec h. \qedhere
\]
\end{proof}

\begin{prop}[reduced polynomial representatives]\label{prop:reduced-rep}
Every class in $W_r(k((x)))/\AS W_r(k((x)))$ has a reduced polynomial representative:
there exists $\vec u\in W_r(k((x)))$ such that $\vec g:=\vf\ominus\AS\vec u$ has all
components in $x^{-1}k[x^{-1}]$ with pole orders either $0$ or prime to $p$.
\end{prop}

\begin{proof}
Induction on $r$.

\emph{Base case $r=1$.} This is the argument of
\texttt{lemma:realization\_over\_Gm}(2) in the thesis: first, writing
$f=f^-+f^+$ with $f^-\in x^{-1}k[x^{-1}]$ and $f^+\in k[[x]]$, surjectivity of $\AS$ on
$k[[x]]$ removes $f^+$; then, as long as the leading term is $cx^{-m}$ with $p\mid m$,
$m>0$, subtracting $\AS(c^{1/p}x^{-m/p})=cx^{-m}-c^{1/p}x^{-m/p}$ strictly decreases the
pole order while staying inside $x^{-1}k[x^{-1}]$; this terminates at $0$ or at pole
order prime to $p$.

\emph{Inductive step.} Let $\vf\in W_r(k((x)))$. By induction there is
$\vec u'\in W_{r-1}(k((x)))$ such that $\vec g':=t(\vf)\ominus\AS\vec u'$ is a reduced
polynomial vector in $W_{r-1}$. Lift $\vec u'$ to $\tilde u:=(\vec u',0)\in
W_r(k((x)))$ and set $\vec h:=\vf\ominus\AS\tilde u$. By
Lemma~\ref{lem:bookkeeping}(i), $t(\vec h)=\vec g'$, i.e.
\[
\vec h=(g'_0,\dots,g'_{r-2},\,h_{r-1}),\qquad h_{r-1}\in k((x)),
\]
with the first $r-1$ components already reduced polynomials. It remains to repair the
last component \emph{without touching the others}; by
Lemma~\ref{lem:bookkeeping}(iii), adjustments of the form
$\ominus\,\AS\bigl(V^{r-1}(z)\bigr)=\ominus\,V^{r-1}(\AS z)$ change only the last
component, by $-\AS(z)$. So:
\begin{itemize}
  \item Write $h_{r-1}=h^-+h^+$ with $h^-\in x^{-1}k[x^{-1}]$, $h^+\in k[[x]]$, choose
        $w\in k[[x]]$ with $\AS w=h^+$ (Lemma~\ref{lem:wp-surj} with $r=1$), and replace
        $\vec h$ by $\vec h\ominus\AS(V^{r-1}w)$: the last component becomes $h^-$.
  \item While the last component has leading term $cx^{-m}$ with $p\mid m$, $m>0$,
        replace $\vec h$ by $\vec h\ominus\AS\bigl(V^{r-1}(c^{1/p}x^{-m/p})\bigr)$: the
        last component changes by
        $-\AS(c^{1/p}x^{-m/p})=-cx^{-m}+c^{1/p}x^{-m/p}$, so its pole order strictly
        decreases and it stays in $x^{-1}k[x^{-1}]$. This terminates at $0$ or at pole
        order prime to $p$.
\end{itemize}
The result is a reduced polynomial vector $\AS$-equivalent to $\vf$.
\end{proof}

\paragraph{The break formula.}
The following is the single external input. It is due, over finite residue fields, to
Kanesaka--Sekiguchi \cite{KanesakaSekiguchi1979} and Thomas \cite{Thomas2005}
(building on Schmid \cite{Schmid1937} for $r=1$), and over arbitrary perfect residue
fields to Elder--Keating \cite{ElderKeating2025}.

\begin{prop}[Schmid--Witt break formula, direction used]\label{prop:break}
Let $K=k((x))$ with $k$ perfect of characteristic $p$, and let $L/K$ be a cyclic totally
ramified extension of degree $p^r$ whose ASW class is represented by a \emph{reduced}
vector $\vec g\in W_r(K)$ (Definition~\ref{def:reduced}), with $m_j:=m(g_j)$. Then the
largest upper-numbering ramification break of $L/K$ equals
\[
u \;=\; \max_{0\le j\le r-1}\; p^{\,r-1-j}\,m_j .
\]
For $r=1$ this is the classical statement (break $=m_0$) already used in the thesis.
\end{prop}

Note that reducedness is exactly what makes the formula exact: two nonzero terms
$p^{r-1-j}m_j=p^{r-1-j'}m_{j'}$ with $j<j'$ would force $p\mid m_{j'}$, so the maximum is
achieved at a unique level and no cancellation of breaks can occur. We use only the
stated direction: \emph{evaluating the formula on one reduced representative computes
$u$}; no minimization over representatives is needed.

\begin{lemma}[realization over $\bG_m$ for $W_r$]\label{lem:realization-Wr}
Let $k$ be algebraically closed of characteristic $p>0$ and let $\cQ$ be a
\emph{connected} $\Z/p^r\Z$-torsor over $\Spec(k((x)))$ with ramification bounded by $d$
(\texttt{definition:local\_ramification\_boundness}). Then there is a reduced polynomial
vector $\vf=(f_0,\dots,f_{r-1})$, $f_j\in x^{-1}k[x^{-1}]$, such that, with
$m_j:=\deg_{x^{-1}}f_j$:
\begin{enumerate}
  \item \label{it:polebound} $p^{\,r-1-j}\,m_j<d$ for all $0\le j\le r-1$; equivalently
        $m_j<d\,p^{\,j-(r-1)}$;
  \item \label{it:torsor} $\cP':=\Spec\Bigl(k[x,x^{-1}][X_0,\dots,X_{r-1}]\big/
        \bigl(\AS\vX\ominus\vf\bigr)\Bigr)$ is a $\Z/p^r\Z$-torsor over $\bG_m$ with
        $\cP'_x\cong\cQ$;
  \item \label{it:infty} $\cP'$ extends to a torsor over $\bP^1_k\setminus\{0\}$ --- in
        particular it is unramified at $\infty$ --- and has ramification bounded by $d$
        at $0$.
\end{enumerate}
\end{lemma}

\begin{proof}
$\cQ=\Spec L$ with $L/K$ cyclic totally ramified of degree $p^r$ and largest upper break
$u<d$. Let $[\vf_0]\in W_r(K)/\AS$ be its ASW class and let $\vf$ be a reduced polynomial
representative (Proposition~\ref{prop:reduced-rep}).

First, $f_0\neq0$: if $f_0=0$ then $\vf=V(\vec b)$ for some $\vec b\in W_{r-1}(K)$, and
$p^{r-1}\vf=(VF)^{r-1}V\vec b=V^{r}F^{r-1}\vec b=0$ in $W_r(K)$, so the class has order
$\le p^{r-1}$, contradicting $[L:K]=p^r$. Hence $m_0\ge1$ and $p\nmid m_0$, and
Proposition~\ref{prop:break} applies:
$\max_j p^{r-1-j}m_j=u<d$, which is \ref{it:polebound}.

For \ref{it:torsor}: by Lemma~\ref{lem:asw-etale} (with $A=k[x,x^{-1}]\ni f_j$), $\cP'$
is a $\Z/p^r\Z$-torsor over $\bG_m$; its restriction to $\Spec k((x))$ is the ASW torsor
of the class of $\vf$, i.e.\ $\cQ$.

For \ref{it:infty}: in the coordinate $y=x^{-1}$ at $\infty$ we have
$f_j\in y\,k[y]$, so $\vf\in W_r(k[y])$ and Lemma~\ref{lem:asw-etale} (with $A=k[y]$)
extends $\cP'$ to a torsor over $\bA^1_k=\Spec k[y]=\bP^1_k\setminus\{0\}$; a torsor
over the full local ring at $\infty$ is unramified there. The ramification bound at $0$
holds because $\cP'_x\cong\cQ$.
\end{proof}

\begin{remark}
Connectedness of $\cQ$ is guaranteed in the application by the standing reduction of
\S1 (inertia equal to all of $G$). This replaces --- and corrects --- the ``take the
principal part'' step: the components $f_j$ are honest polynomials in $x^{-1}$, so
substituting $x\mapsto x_i$ and forming the Witt sum below are purely algebraic
operations, and all subsequent $\AS$-adjustments were already performed \emph{before}
symmetrization, inside $W_r(k((x)))$.
\end{remark}

\begin{remark}[sanity check]\label{rem:sanity}
For every $j$ with $p^{\,r-1-j}\ge d$, the bound \ref{it:polebound} forces $m_j=0$, i.e.\
$f_j=0$: all sufficiently low Witt components of a reduced representative vanish
outright. This makes visible how strong the hypothesis $u<d$ is at the lower levels.
\end{remark}

%====================================================================
\section{Gap 2: the symmetric-power cover at $\eta$ is the ASW cover of the Witt sum}
%====================================================================

We now generalize the invariant-ring computation of the thesis's $r=1$ proof
(\texttt{theorem:ramification\_after\_blowup\_P1}, Artin--Schreier case) to Witt
coordinates. Fix $\vf\in W_r(R)$, $R=k[x,x^{-1}]$, and let
\[
S \;=\; R[\vX]\big/\bigl(\AS\vX\ominus\vf\bigr)
\]
be the coordinate ring of the torsor $\cP'\to\bG_m$ of
Lemma~\ref{lem:realization-Wr}. Write $R^{\otimes_kd}=k[x_1^{\pm1},\dots,x_d^{\pm1}]$
and
\[
S^{\otimes_kd}
= R^{\otimes_kd}[\vX_1,\dots,\vX_d]\big/\bigl(\AS\vX_1\ominus\vf(x_1),\dots,
\AS\vX_d\ominus\vf(x_d)\bigr),
\]
where $\vX_i=(X_{i,0},\dots,X_{i,r-1})$ and $\vf(x_i)$ is the image of $\vf$ under
$x\mapsto x_i$. Exactly as in the thesis, the $d$-fold product torsor
$p_1^{-1}\cP'\otimes\cdots\otimes p_d^{-1}\cP'$ on $U^d$ is by definition the quotient of
$\Spec S^{\otimes_kd}$ by $G^{d-1}$, embedded in $G^d$ as the kernel of the sum map, and
acting (in Witt coordinates, writing $\vec g$ for the image of $g\in G$ under
$\Z/p^r\Z\cong W_r(\F_p)$) by
\[
(g_1,\dots,g_{d-1})\colon\quad
\vX_1\mapsto\vX_1\oplus\vec g_1,\qquad
\vX_i\mapsto\vX_i\oplus\vec g_i\ominus\vec g_{i-1}\ (1<i<d),\qquad
\vX_d\mapsto\vX_d\ominus\vec g_{d-1}.
\]

\begin{prop}[identification of the product torsor]\label{prop:cover-identification}
Set
\[
\vY:=\bigoplus_{i=1}^d\vX_i\in W_r\bigl(S^{\otimes_kd}\bigr),
\qquad
\vec\alpha:=\bigoplus_{i=1}^d\vf(x_i)\in W_r\bigl(R^{\otimes_kd}\bigr),
\]
where $\bigoplus$ denotes the $d$-fold Witt vector sum $\oplus$ of
\eqref{eq:addition-shape} (not a direct sum); in particular $\vY$ is a single Witt
vector of length $r$. Then:
\begin{enumerate}
  \item \label{it:inv} Every component $Y_n$ of $\vY$ is $G^{d-1}$-invariant, and
        \[
        \bigl(S^{\otimes_kd}\bigr)^{G^{d-1}}
        \;\cong\;
        R^{\otimes_kd}[\vY]\big/\bigl(\AS\vY\ominus\vec\alpha\bigr)
        \]
        as $G$-torsors over $R^{\otimes_kd}$ (the residual $G=G^d/G^{d-1}$ acting by
        $\vY\mapsto\vY\oplus\vec g$).
  \item \label{it:symm} Every component $\alpha_n$ of $\vec\alpha$ is a symmetric Laurent
        polynomial, i.e.\ lies in
        $\bigl(R^{\otimes_kd}\bigr)^{S_d}=k[e_1,\dots,e_d,e_d^{-1}]$, and
        \[
        \Bigl(\bigl(S^{\otimes_kd}\bigr)^{G^{d-1}}\Bigr)^{S_d}
        \;\cong\;
        k[e_1,\dots,e_d,e_d^{-1}][\vY]\big/\bigl(\AS\vY\ominus\vec\alpha\bigr).
        \]
\end{enumerate}
Consequently, restricting to the complete local field
$\widehat K_\eta=\kappa(\eta)((e_d))$, $\kappa(\eta)=k(e_1/e_d,\dots,e_{d-1}/e_d)$, of
\textup{\texttt{eq:complete\_field\_at\_generic\_point}},
\[
\cP'^{[d]}\big|_{\Spec\widehat K_\eta}
\;\cong\;
\Spec\;\widehat K_\eta[\vZ]\big/\bigl(\AS\vZ\ominus\vec\alpha\bigr).
\]
\end{prop}

\begin{proof}
\ref{it:inv} An element $\sigma=(g_1,\dots,g_{d-1})\in G^{d-1}$ acts as a ring
automorphism of $S^{\otimes_kd}$: this is its \emph{original} action, defined on the
generators $X_{i,n}$ by the twisted formulas above (which mix components through the
carry polynomials of \eqref{eq:addition-shape}). By functoriality of $W_r$ (see the
conventions), the induced map $W_r(\sigma)$ on $W_r(S^{\otimes_kd})$ is componentwise
and commutes with $\oplus,\ominus$; on the generators it reproduces the twisted action,
$W_r(\sigma)(\vX_i)=\vX_i\oplus\vec g_i\ominus\vec g_{i-1}$ (with $\vec g_0=\vec
g_d=\vec 0$). Hence, using commutativity and associativity of $\oplus$, the Witt sum
transforms by the telescoping constant
\[
W_r(\sigma)(\vY)\;=\;\bigoplus_{i=1}^dW_r(\sigma)(\vX_i)\;=\;\vY\oplus
\vec g_1\oplus(\vec g_2\ominus\vec g_1)\oplus\cdots\oplus(\ominus\,\vec g_{d-1})
=\vY.
\]
Since $W_r(\sigma)$ is componentwise, this equality of $r$-tuples says exactly
$\sigma(Y_n)=Y_n$ for every $n$, i.e.\ each $Y_n\in S^{\otimes_kd}$ is invariant. Since $\AS=F-\mathrm{id}$ is additive
($F$ is a ring endomorphism),
\[
\AS\vY=\bigoplus_{i=1}^d\AS\vX_i=\bigoplus_{i=1}^d\vf(x_i)=\vec\alpha .
\]
Hence there is a well-defined $R^{\otimes_kd}$-algebra homomorphism
\[
\phi\colon\; T:=R^{\otimes_kd}[\vZ]\big/\bigl(\AS\vZ\ominus\vec\alpha\bigr)
\;\longrightarrow\;\bigl(S^{\otimes_kd}\bigr)^{G^{d-1}},
\qquad \vZ\mapsto\vY .
\]
By Lemma~\ref{lem:asw-etale}, $\Spec T$ is a $G$-torsor over $U^d$. The target is the
product torsor (the contracted product of the $p_i^{-1}\cP'$ along the sum map
$G^d\to G$), a $G$-torsor over $U^d$ by the thesis's construction of torsor products;
its residual $G$-action is induced by $(g,0,\dots,0)\in G^d$, i.e.\
$\vX_1\mapsto\vX_1\oplus\vec g$, under which $\vY\mapsto\vY\oplus\vec g$. So
$\Spec\phi$ is a $G$-equivariant morphism of $G$-torsors over $U^d$, hence an
isomorphism by Lemma~\ref{lem:torsor-mor}.

\ref{it:symm} $S_d$ acts on $S^{\otimes_kd}$ by permuting the tensor factors, i.e.\
$x_i\mapsto x_{\sigma(i)}$, $\vX_i\mapsto\vX_{\sigma(i)}$. As in \ref{it:inv}, each
$\sigma\in S_d$ is a ring automorphism, so $W_r(\sigma)$ is componentwise and commutes
with $\oplus$. Since $\oplus$ is commutative
and associative, permuting the summands of a Witt sum leaves it unchanged; thus
$\vY$ is $S_d$-invariant (componentwise, as before) and
$W_r(\sigma)(\vec\alpha)=\bigoplus_i\vf(x_{\sigma(i)})=\vec\alpha$, so each
$\alpha_n\in R^{\otimes_kd}$ is a symmetric Laurent polynomial, hence lies in
$k[e_1,\dots,e_d,e_d^{-1}]$ (the $S_d$-invariants of $R^{\otimes_kd}$, as in the
thesis). By Lemma~\ref{lem:asw-etale}, $T$ is free over $R^{\otimes_kd}$ with basis the
monomials $\prod_nY_n^{a_n}$, $0\le a_n<p$, each of which is $S_d$-invariant; writing an
element in this basis, it is $S_d$-invariant if and only if all its coefficients are.
Therefore
\[
T^{S_d}=\bigl(R^{\otimes_kd}\bigr)^{S_d}[\vY]\big/\bigl(\AS\vY\ominus\vec\alpha\bigr)
=k[e_1,\dots,e_d,e_d^{-1}][\vY]\big/\bigl(\AS\vY\ominus\vec\alpha\bigr).
\]

The final statement follows exactly as in the $r=1$ case: by definition $\cP'^{[d]}$ is
the $\Xi=G^{d-1}\rtimes S_d$-quotient above, and its restriction to
$\Spec\widehat K_\eta$ along
$k[e_1,\dots,e_d,e_d^{-1}]\to\widehat K_\eta$ is the ASW cover of the image of
$\vec\alpha$ in $W_r(\widehat K_\eta)$.
\end{proof}

\begin{remark}
The proof via Lemma~\ref{lem:torsor-mor} replaces the rank count of the $r=1$ argument
($p^{rd}/p^{r(d-1)}=p^r$ together with irreducibility of the invariant's minimal
equation); the rank count alone would still require showing that $\vY$ generates the
invariant ring, which is exactly what equivariance gives for free. Note also that a pure
restriction argument on $H^1$'s (``the class of the product torsor is the sum of the
pulled-back characters'') identifies the cover over $k(x_1,\dots,x_d)$ but does
\emph{not} by itself descend it to the symmetric base $k(e_1,\dots,e_d)$; the
invariant-theoretic computation above is what effects the descent.
\end{remark}

%====================================================================
\section{The degree count}
%====================================================================

This section reproduces, for completeness, the two elementary lemmas and the estimate of
the earlier note (verified line by line in the review), now fed with the legitimate
inputs provided by Lemma~\ref{lem:realization-Wr} and
Proposition~\ref{prop:cover-identification}. Throughout, $\vf$ is the reduced polynomial
vector of Lemma~\ref{lem:realization-Wr}, $m_j=\deg_{x^{-1}}f_j$ satisfies
$m_j<d\,p^{\,j-(r-1)}$, and $y_i:=1/x_i$, so that $f_j(x_i)\in y_i\,k[y_i]$ is a
polynomial in $y_i$ of degree $m_j$.

\begin{lemma}[symmetric regularity]\label{lem:symreg}
Let $\Phi\in k[y_1,\dots,y_d]^{S_d}$ be symmetric with every monomial of total degree
$<d$. Then $\Phi\in\kappa(\eta)=k(e_1/e_d,\dots,e_{d-1}/e_d)$; in particular $\Phi$ is
regular at $\eta$.
\end{lemma}

\begin{proof}
$\Phi$ is a $k$-combination of monomial symmetric functions $m_\lambda(y)$ with
$|\lambda|<d$. Each $m_\lambda(y)$ is an isobaric polynomial in the elementary symmetric
functions $e_\mu(y)=\prod_ie_{\mu_i}(y)$ with $\mu\vdash|\lambda|$; every part
$\mu_i\le|\lambda|<d$, so only $e_1(y),\dots,e_{d-1}(y)$ occur --- never $e_d(y)$. Now
$e_k(y)=e_{d-k}(x)/e_d(x)$, which for $1\le k\le d-1$ lies in $\kappa(\eta)$. Hence
$\Phi\in\kappa(\eta)\subset\widehat\Oo_\eta=\kappa(\eta)[[e_d]]$.
\end{proof}

\begin{lemma}[isobaricity of Witt addition]\label{lem:isobaric}
Give the variable $f_j(x_i)$ the weight $p^j$. Then each component
$\alpha_n=\bigl(\bigoplus_{i=1}^d\vf(x_i)\bigr)_n$ is isobaric of weight $p^n$: every
monomial $\prod_{i,j}\bigl(f_j(x_i)\bigr)^{e_{ij}}$ of $\alpha_n$ satisfies
$\sum_{i,j}e_{ij}\,p^{\,j}=p^{\,n}$.
\end{lemma}

\begin{proof}
The $n$-th ghost component $w_n(\va)=\sum_{j\le n}p^ja_j^{p^{n-j}}$ is isobaric of weight
$p^n$ when $\deg a_j=p^j$. The universal addition polynomials $S_n$ are determined over
$\Z$ by additivity of the ghost components and the recursion
$p^nS_n=\sum_jp^j\bigl(a_j^{p^{n-j}}+b_j^{p^{n-j}}\bigr)-\sum_{j<n}p^jS_j^{p^{n-j}}$,
which preserves isobaricity of weight $p^n$ by induction; the same holds after reduction
mod $p$ and under $d$-fold iteration of $\oplus$. Hence $\alpha_n$, a universal isobaric
polynomial evaluated at the $f_j(x_i)$, is isobaric of weight $p^n$.
\end{proof}

\begin{theorem}[regularity of all Witt components]\label{thm:regularity}
With $\vf$ as in Lemma~\ref{lem:realization-Wr}, every component $\alpha_n$
$(0\le n\le r-1)$ of $\vec\alpha=\bigoplus_{i=1}^d\vf(x_i)$ lies in $\kappa(\eta)$; in
particular $\vec\alpha\in W_r\bigl(\widehat\Oo_\eta\bigr)$.
\end{theorem}

\begin{proof}
Fix $n\le r-1$ and a monomial $\prod_{i,j}\bigl(f_j(x_i)\bigr)^{e_{ij}}$ of $\alpha_n$;
set $E_j:=\sum_ie_{ij}$, so $\sum_jE_jp^{\,j}=p^{\,n}$ by Lemma~\ref{lem:isobaric}.
Expanding each $f_j(x_i)$ as a polynomial in $y_i$ of degree $m_j$, every resulting
monomial in the $y_i$ has total degree
\[
\le\;\sum_{i,j}e_{ij}\,m_j\;=\;\sum_jE_j\,m_j
\;<\;\sum_jE_j\,d\,p^{\,j-(r-1)}
\;=\;d\,p^{-(r-1)}\sum_jE_jp^{\,j}
\;=\;d\,p^{\,n-(r-1)}\;\le\;d,
\]
strictly: some $E_j>0$ (as $\sum_jE_jp^j=p^n>0$), and the pole bound
$m_j<d\,p^{\,j-(r-1)}$ of Lemma~\ref{lem:realization-Wr}\ref{it:polebound} is strict
even when $m_j=0$. Thus every monomial of $\alpha_n$ --- a symmetric polynomial in
$y_1,\dots,y_d$ by Proposition~\ref{prop:cover-identification}\ref{it:symm} --- has
total $y$-degree $<d$, and Lemma~\ref{lem:symreg} gives $\alpha_n\in\kappa(\eta)$.
\end{proof}

%====================================================================
\section{Proof of the theorem}
%====================================================================

\begin{proof}[Proof of Theorem~\ref{thm:main}]
By the standing reduction of \S1 we may assume $k$ algebraically closed and $\cP\to U$
totally ramified at $P$ with inertia all of $G$, so that $\cQ:=\cP_x$ is a connected
$\Z/p^r\Z$-torsor over $\Spec k((x))$ with ramification bounded by $d=\deg\ndls$.

Lemma~\ref{lem:realization-Wr} produces a torsor $\cP'\to\bG_m$, defined by a reduced
polynomial vector $\vf$ with $p^{r-1-j}m_j<d$, with $\cP'_x\cong\cQ$. By the reduction
to $\bP^1$ (\texttt{cor:reduction\_to\_P1}, valid for arbitrary finite abelian $G$), the
ramification of $\cP^{[d]}$ at $\eta_\ndls$ agrees with that of $\cP'^{[d]}$ at
$\eta_{d\cdot0}$; so we may assume $C=\bP^1_k$, $\mdls=d\cdot0+\infty$,
$\ndls=d\cdot0$, $\cP=\cP'$.

By Proposition~\ref{prop:cover-identification},
\[
\cP'^{[d]}\big|_{\Spec\widehat K_\eta}
\;\cong\;
\Spec\;\widehat K_\eta[\vZ]\big/\bigl(\AS\vZ\ominus\vec\alpha\bigr),
\qquad
\vec\alpha=\bigoplus_{i=1}^d\vf(x_i),
\]
with $\widehat K_\eta=\kappa(\eta)((e_d))$ the complete local field at $\eta_\ndls$ and
$\widehat\Oo_\eta=\kappa(\eta)[[e_d]]$ its valuation ring.

By Theorem~\ref{thm:regularity}, $\vec\alpha\in W_r(\kappa(\eta))\subset
W_r(\widehat\Oo_\eta)$. Hence, by Lemma~\ref{lem:asw-etale} applied over
$A=\widehat\Oo_\eta$, the ASW cover $\AS\vZ=\vec\alpha$ is finite \'etale over
$\widehat\Oo_\eta$, and $\cP'^{[d]}\big|_{\Spec\widehat K_\eta}$ is its generic fibre:
every field factor of $\widehat K_\eta[\vZ]/(\AS\vZ\ominus\vec\alpha)$ is the fraction
field of a finite \'etale $\widehat\Oo_\eta$-algebra, i.e.\ an unramified extension of
$\widehat K_\eta$ with respect to the discrete valuation ring $\widehat\Oo_\eta$.
By \texttt{definition:tamely\_ramified\_in\_codim\_1}, $\cP^{[d]}$ is unramified at
$\eta_\ndls$.
\end{proof}

\begin{remark}[this is a direct proof, not an induction]\label{rem:direct}
The argument proves all $r$ simultaneously; it reproves the case $r=1$
(\texttt{cor:ramification\_after\_blowup\_is\_tame}(1)) as the special case of a vector
of length one and never invokes it. In particular, when merging into the thesis, the
current proof body of
\texttt{theorem:ramification\_after\_blowup\_is\_tame\_for\_p\_power} from the sentence
``We proceed by induction on $r$'' onward should be \emph{replaced} by
\S\S2--5 above --- including the step which applies the $r=1$ theorem to the top layer
$\cP\to\cP'$: that step is unjustified, since the hypothesis bounds the ramification of
$\cP$ in \emph{upper} numbering, while the break of the top degree-$p$ layer is the
largest \emph{lower}-numbering break, which can be of size about $p^{r-1}d$. The
reduction paragraph preceding it (algebraically closed base, totally ramified, inertia
all of $G$) is correct for all $r$ and is exactly the standing reduction of \S1.
\end{remark}

\begin{remark}[bookkeeping of the estimate]
As in the earlier note, the bound is exact: the weight $p^{\,j-(r-1)}$ is paired against
the isobaricity constraint $\sum_jE_jp^{\,j}=p^{\,n}$, so the estimate
$d\,p^{\,n-(r-1)}\le d$ is independent of how the Witt carries distribute the exponents
$E_j$; the hypothesis ``break $<d$ at $p_\infty$'' is used at \emph{every} Witt level at
once. Remark~\ref{rem:sanity} records the extreme case: components with
$p^{\,r-1-j}\ge d$ vanish outright in the reduced representative.
\end{remark}

\begin{thebibliography}{9}

\bibitem{Schmid1937}
H.~L.~Schmid,
\emph{Zur Arithmetik der zyklischen $p$-K\"orper},
J. Reine Angew. Math. \textbf{176} (1937), 161--167.

\bibitem{KanesakaSekiguchi1979}
K.~Kanesaka and K.~Sekiguchi,
\emph{Representation of Witt vectors by formal power series and its applications},
Tokyo J. Math. \textbf{2} (1979), no.~2, 349--370.

\bibitem{Thomas2005}
L.~Thomas,
\emph{Ramification groups in Artin--Schreier--Witt extensions},
J. Th\'eor. Nombres Bordeaux \textbf{17} (2005), no.~2, 689--720.

\bibitem{ElderKeating2025}
G.~G.~Elder and K.~Keating,
\emph{Artin--Schreier--Witt extensions and ramification breaks},
preprint (2025), arXiv:2503.16830.

\end{thebibliography}

\end{document}
